\section{Conclusion}

This work presented a complete pipeline for content- and preference-aware music recommendation grounded in a structured music knowledge graph. We constructed two complementary RDF graph representations of 7,013 tracks with verified ontological alignment to established vocabularies and Wikidata, validated through SPARQL queries in GraphDB. Genre ranking comparisons between the simple and rich KG variants illustrated the expressive value of weighted representations, while decade temporal chains and Wikidata super-class alignment confirmed structural correctness throughout.

A symbolic-feature autoencoder successfully compressed 1,525-dimensional jSymbolic/music21 descriptors to a 128-dimensional bottleneck (final MSE: 0.62), and KG embeddings provided structural node initialisation for all entity types. Baseline evaluation confirms that KNN collaborative filtering substantially outperforms popularity-based recommendation in both ranking quality and catalogue diversity (Overall Score 0.420 vs.\ 0.134), while the XGBoost hybrid reveals the retrieval bottleneck in content-only approaches without graph structure.

Full training and evaluation of the HGT recommender with the popularity-debiased listwise loss remains ongoing; computational constraints limited the number of experiments that could be completed within the project timeline. Future work will complete the HGT evaluation and compare its performance against the presented baselines, examine the effect of KGE initialisation strategies on recommendation quality, and investigate how Wikidata-grounded graph semantics influence recommendation diversity across musical dimensions.

A structured hyperparameter search---covering the number of HGT attention heads, hidden dimensionality, temperature~$\tau$ and debiasing strength~$\lambda$ of the listwise loss, and the KGE embedding dimension---would help identify the best-performing configuration and reduce the sensitivity to ad-hoc choices made under time pressure. Complementing this, a rigorous ablation study isolating the individual contributions of KG structure, symbolic audio features, and the popularity-debiased loss would clarify which components drive performance. Extending the baseline suite to include graph-collaborative methods such as LightGCN~\cite{he2020lightgcn} or KGCN~\cite{wang2019kgcn}, combined with paired statistical significance tests (e.g.\ Wilcoxon signed-rank on per-user NDCG), would then provide the evidence needed to determine whether the additional complexity of heterogeneous graph modelling is justified for this research question.
